\subsection{The complete peripheral fibre in a universal coprime Frattini cover}

\begin{statement}{D.39}{Proved}{Actual automorphisms and every generating lift at the first coprime characteristic cover}
Let \(q:\widehat F\twoheadrightarrow D\) be an epimorphism to a
finite group, let \(K=\ker q\), and let \(p\nmid |D|\).
Put
\[
 \Phi_p(K)=\overline{K^p[K,K]},\qquad
 E=\widehat F/\Phi_p(K),\qquad
 M=K/\Phi_p(K),\qquad n=|D|.
 \tag{D.39.1}
\]
Then \(E\) is finite of order \(n p^{n+1}\), and \(M\) is the
actual \(D\)-module of first homology of the regular pants cover.
If \(K\) is characteristic in \(\widehat F\), so is
\(\Phi_p(K)\).

Write \(X_i\) for the images of the three peripheral generators
in \(E\), \(d_i=q(x_i)\), and \(e_i=\operatorname{ord}(d_i)\).
For every \(\alpha\in\operatorname{Aut}(D)\),
\(\lambda\in\widehat{\mathbf Z}^{\times}\), and specified
\(h_i\in D\) with
\[
                 \alpha(d_i)=h_i d_i^\lambda h_i^{-1},
 \tag{D.39.2}
\]
there is an actual automorphism \(\widetilde\alpha\) of \(E\)
covering \(\alpha\) such that, for any chosen lifts \(H_i\) of
\(h_i\),
\[
 \widetilde\alpha(X_i)
       =a_iH_iX_i^\lambda H_i^{-1}a_i^{-1}
                     \quad\text{with }a_i\in M.
 \tag{D.39.3}
\]
Every generating source lift above \(\alpha q\) factors
through an automorphism of this same \(E\).

Let \(B\subset M\) be the boundary submodule and \(C=M/B\).
The complete fibre with the base conjugators \(h_i\) fixed
has, after stated linear choices, the parameter set
\[
 (M/M^D)\ \times\
 \operatorname{Hom}_{\mathbf F_p[D]}(C,B)\ \times\
 \operatorname{Aut}_{\mathbf F_p[D]}(C).
 \tag{D.39.4}
\]
Its literal cardinality is
\[
 p^{n-1}
 \bigl|\operatorname{Hom}_{\mathbf F_p[D]}(C,B)\bigr|
 \bigl|\operatorname{Aut}_{\mathbf F_p[D]}(C)\bigr|.
 \tag{D.39.5}
\]
Allowing every ordinary base conjugator multiplies this
cardinality by
\[
                       \prod_{i=1}^3
                       \frac{|C_D(d_i)|}{e_i}.
 \tag{D.39.6}
\]
These formulas count literal generating source maps, or
equivalently actual automorphisms of \(E\), before taking
any outer quotient.
In particular all D.37 parameters have the explicitly
described nonempty full fibre at each first universal
elementary-\(p\) cover with \(p\ne5\).
\end{statement}

\begin{proof}
Use the left \(D\)-modules of the right Cayley graph of
\(D\) with its two directed generator edges:
\[
 C_1=\mathbf F_p[D]e_x\oplus\mathbf F_p[D]e_y,\qquad
 C_0=\mathbf F_p[D],\qquad
 \partial(de_x)=d d_1-d,\quad
 \partial(de_y)=d d_2-d.
 \tag{D.39.7}
\]
Every based word \(w\) has its literal edge chain
\(v(w)\), with
\[
 v(uv)=v(u)+q(u)v(v),\qquad
 \partial v(w)=q(w)-1.
 \tag{D.39.8}
\]
Consequently
\[
 w\longmapsto(q(w),v(w)),\qquad
 (d,v)(e,w)=(de,v+dw)
 \tag{D.39.9}
\]
is an actual finite quotient map of \(\widehat F\).
Choose a spanning tree and based words \(t_d\) for its
vertices.  The non-tree Schreier words
\(t_d x_jt_{d d_j}^{-1}\) give a free basis of the
finite-index subgroup of the discrete free group.
Their \(n+1\) edge cycles form a basis of \(\ker\partial\).
The same words freely generate the open free profinite
subgroup \(K\), by the finite-coset Schreier rewriting
and its inverse on the two generators.
Thus their map to \(\ker\partial\) has kernel exactly
\(\Phi_p(K)\).  This proves
\[
 M=\ker\partial,\qquad
 0\longrightarrow M\longrightarrow C_1
   \xrightarrow{\partial} C_0
   \xrightarrow{\mathrm{aug}}\mathbf F_p\longrightarrow0
 \tag{D.39.10}
\]
with the actual source map \textup{(D.39.9)}, and proves
the asserted finiteness and order.

Here is an explicit splitting of this extension.
Put \(b_d=v(t_d)\) and
\[
 \kappa=n^{-1}\sum_{d\in D}b_d\in C_1,\qquad
 s(d)=(d,\kappa-d\kappa).
 \tag{D.39.11}
\]
The vector \(\kappa-d\kappa-b_d\) belongs to
\(M\), since its boundary is zero.  Hence \(s(d)\)
belongs to the actual image in \textup{(D.39.9)}.
Direct multiplication gives \(s(d)s(e)=s(de)\).
In these coordinates the exact isomorphism
\[
 E\longrightarrow M\rtimes D,\qquad
 (d,v)\longmapsto(v-\kappa+d\kappa,d)
 \tag{D.39.12}
\]
has inverse \((m,d)\mapsto(d,m+\kappa-d\kappa)\).
This realizes the splitting without a change of the
source peripheral elements.

We next retain the complete boundary module.
For \(i=1,2,3\), let
\[
 b_i=v(x_i^{e_i})=X_i^{e_i}\in M.
 \tag{D.39.13}
\]
In the second equality \(M\) is written additively.
The cycles \(d b_i\) are indexed by the right cosets
\(d\langle d_i\rangle\), with the left deck action.
Filling the boundary components of the connected
oriented pants cover gives a closed connected
oriented surface.  Its cellular two-cycles have
rank one, generated by the oriented fundamental
class.  Therefore the entire space of relations
among these boundary cycles is exactly their sum:
\[
 0\longrightarrow\mathbf F_p
 \xrightarrow{\,1\mapsto\Delta\,}
 \bigoplus_{i=1}^3
       \mathbf F_p[D/\langle d_i\rangle]
 \longrightarrow B\longrightarrow0,\qquad
 \Delta=\sum_{i,d\langle d_i\rangle}
                    [d\langle d_i\rangle].
 \tag{D.39.14}
\]
One may also obtain this sequence directly by attaching
one two-cell to each literal boundary cycle in
\textup{(D.39.7)}; each directed edge then occurs with
both opposite induced orientations.
The quotient \(C=M/B\) is the homology of that
compact surface.  In particular, with
\(b=\sum_i n/e_i\),
\[
 \dim M=n+1,\quad \dim B=b-1,\quad
 \dim C=n+2-b.
 \tag{D.39.15}
\]

Let \(\bar\lambda\in\mathbf F_p^\times\) be the
reduction of \(\lambda\).  The prescribed map on
the boundary cycles is
\[
 T_B(d b_i)=\bar\lambda\,\alpha(d)h_i b_i.
 \tag{D.39.16}
\]
It is well defined: replacing \(d\) by \(d d_i^j\)
changes the right side by
\(\alpha(d_i)^j h_i=h_i d_i^{\lambda j}\),
which fixes \(b_i\).
For each \(i\), the map
\[
 d\langle d_i\rangle\longmapsto
            \alpha(d)h_i\langle d_i\rangle
 \tag{D.39.17}
\]
is a bijection, with inverse obtained from
\(\alpha^{-1}\) and \(h_i^{-1}\).
Thus it sends \(\Delta\) to
\(\bar\lambda\Delta\) and induces an invertible
\(\alpha\)-semilinear map \(T_B:B\to B\).
All boundary path choices in this step are
specified by the actual coset map
\textup{(D.39.17)}.

Because \(p\nmid n\), averaging any projection
over \(D\) gives a \(D\)-linear projection.
Every finite \(\mathbf F_p[D]\)-module is
therefore semisimple.
Splitting \textup{(D.39.10)} gives, by cancellation
of simple summands,
\[
                    M\simeq\mathbf F_p\oplus\mathbf F_p[D].
 \tag{D.39.18}
\]
The regular-module isomorphism
\(d\mapsto\alpha(d)\) shows that its
\(\alpha\)-twist has the same simple multiplicities.
Since \(T_B\) already identifies the two boundary
summands, cancellation identifies their complementary
summands as well.
It follows that \(T_B\) extends to an invertible
\(\alpha\)-semilinear map
\[
              T:M\longrightarrow M,\qquad
              T(dm)=\alpha(d)T(m).
 \tag{D.39.19}
\]
This is an effective finite construction: form the
matrices \textup{(D.39.7)--(D.39.17)}, solve the
linear intertwining and boundary equations, and
enumerate their finite solution space until a
matrix of nonzero determinant is found.
The preceding semisimple argument proves that
this enumeration terminates.

The actual group automorphism is now
\[
 (m,d)\longmapsto(Tm,\alpha(d))
             \quad\text{in \textup{(D.39.12)}}.
 \tag{D.39.20}
\]
Its inverse uses \(T^{-1}\) and \(\alpha^{-1}\).
Set \(U_i=H_iX_i^\lambda H_i^{-1}\) and
\(Y_i=\widetilde\alpha(X_i)\).
Their projections agree by \textup{(D.39.2)}, and
\[
 Y_i^{e_i}=T b_i=\bar\lambda h_i b_i=U_i^{e_i}
 \tag{D.39.21}
\]
by \textup{(D.39.16)}.
Write \(Y_i=v_iU_i\), \(v_i\in M\), and let
\(A_i\) be the action of \(\alpha(d_i)\) on \(M\).
Equation \textup{(D.39.21)} says
\((1+A_i+\cdots+A_i^{e_i-1})v_i=0\).
The explicit kernel conjugator is
\[
       a_i=-e_i^{-1}\sum_{j=0}^{e_i-1}jA_i^j v_i.
 \tag{D.39.22}
\]
Indeed
\((1-A_i)\sum_j jA_i^j
 =(1+\cdots+A_i^{e_i-1})-e_i\),
so \((1-A_i)a_i=v_i\).
Consequently \(a_iU_i a_i^{-1}=Y_i\), which proves
\textup{(D.39.3)} with every chosen \(H_i\) retained.

We now prove that these constructions parameterize
the complete generating fibre.
Any homomorphism \(f:\widehat F\to E\) with
\(\pi f=\alpha q\) maps \(K\) into the elementary
abelian group \(M\).  It therefore kills
\(\overline{K^p[K,K]}\) and factors through the
actual quotient in \textup{(D.39.1)}.
It is onto precisely when the resulting
endomorphism of finite \(E\) is an automorphism.
Such an automorphism restricts to an invertible
\(\alpha\)-semilinear \(T\) on \(M\).
In \textup{(D.39.12)} its value on \(s(d)\)
has the form \((z(d),\alpha(d))\), where
\[
 z(de)=z(d)+\alpha(d)z(e).
 \tag{D.39.23}
\]
Summing over \(e\) gives
\(z(d)=v-\alpha(d)v\), with
\(v=n^{-1}\sum_e z(e)\).
Thus every automorphism is uniquely
\(\operatorname{Inn}(v)\) followed by
\textup{(D.39.20)}, with \(v\) taken modulo
\(M^D\).
Raising the three peripheral identities to
their \(e_i\)th powers shows that the
fixed-conjugator condition is exactly
\(T|_B=T_B\); its sufficiency was proved in
\textup{(D.39.21)--(D.39.22)}.

Choose a \(D\)-linear complement \(C_0\) to
\(B\) in \(M\), identify it with \(C=M/B\),
and choose one extension \(T_0\) in
\textup{(D.39.19)}.
Every other extension is uniquely \(T_0 R\), where
in this chosen direct sum
\[
       R=\begin{pmatrix}1&u\\0&A\end{pmatrix},
       \qquad
       u\in\operatorname{Hom}_{\mathbf F_p[D]}(C,B),
       \quad A\in\operatorname{Aut}_{\mathbf F_p[D]}(C).
 \tag{D.39.24}
\]
This proves the actual parameterization
\textup{(D.39.4)}.
By \textup{(D.39.18)}, \(\dim M^D=2\), so
\(|M/M^D|=p^{n-1}\), which proves
\textup{(D.39.5)}.

Finally every ordinary base conjugator is
\(h_i c_i\), with \(c_i\in C_D(d_i)\).
The new boundary map is determined by the
three vectors \(c_i b_i\).
The stabilizer in \(D\) of \(b_i\) is exactly
\(\langle d_i\rangle\).
To verify this in the actual module, use
\textup{(D.39.14)}: two distinct individual
boundary basis vectors cannot become equal
modulo the one all-boundary sum, since there
are at least three boundary components.
Thus the distinct maps on the \(i\)th family
number \(|C_D(d_i)|/e_i\).
The three families may be chosen independently;
their sole relation, \(\Delta\), is preserved
by every such choice.
Distinct resulting \(T_B\)'s give disjoint
automorphism fibres, each of cardinality
\textup{(D.39.5)}.  This proves the complete
ordinary union formula \textup{(D.39.6)}.

If \(D\) is a five-group and \(p\ne5\),
the projection \(E\to D\) is its maximal
pro-five quotient: the normal \(p\)-kernel
has trivial image in every finite five-group.
For \(D=P_n\), D.37 supplies the characteristic
kernel, and characteristicness of
\(\Phi_p(K)\) follows from its displayed
definition.  This proves the fixed-source
application stated above.
\end{proof}

The program \texttt{a\_coprime\_frattini.py} constructs
the actual Cayley chains, all Schreier basis words,
the boundary sum and its full rank, the splitting
\textup{(D.39.11)}, the semilinear matrix and its
inverse, and all three conjugators
\textup{(D.39.22)}.
It checks 28 cases over \(C_5\) and \(C_5^2\)
in characteristics two and three, and over \(S_3\)
in characteristic five, including all twelve
base-centralizer triples in the latter case.
The separate
\texttt{a\_coprime\_frattini\_count.py}
enumerates the entire source fibre in actual
groups of orders 48 and 54.
For the genus-one \(C_3\)-cover in characteristic
two it finds sixteen weak maps and twelve
epimorphisms; for the genus-zero \(C_3\)-cover
it finds twelve maps, all onto, with three
distinct ordinary boundary maps.
For the \(C_2\)-cover in characteristic three
it finds six maps, all onto, with two distinct
ordinary boundary maps.
Every multiplier tested agrees with
\textup{(D.39.5)--(D.39.6)}.
All source files, matrices and exact outputs
are retained in \texttt{code/}.
